%================================================================
\section{History matching and the image}\label{sec:hm}

\subsection{Ruling points out}

History matching asks which parameter points the targets cannot rule out. Let
$\widehat M(c)$ be the simulated moments at $c$. The implausibility of $c$ is
\begin{equation}\label{eq:implaus}
  I(c)
  \;=\;
  \bigl(z-\widehat M(c)\bigr)^{\!\top}
  V(c)^{-1}
  \bigl(z-\widehat M(c)\bigr),
  \qquad
  V(c)=\mathrm{diag}\bigl(\sigma^2_z+\sigma^2_{\mathrm{MC}}(c)+\delta^{\circ 2}\bigr),
\end{equation}
a variance-weighted sum of squared gaps between targets and simulated moments.
The three terms of $V$ are the Monte Carlo variance of the targets, the Monte
Carlo variance of the candidate's moments, and a fixed discrepancy $\delta$, the
allowance for a model that does not match the world exactly
\citep{kennedy2001,vernon2010}. In this planted exercise $\delta$ is a design
tolerance rather than an estimated model-error variance:
$\delta=(\DeltaOne,\,\DeltaTwo,\,\DeltaThree,\,\DeltaFour)$ in the order of
Table~\ref{tab:moments}, that is, $\DeltaOne$ for the two levels and $\DeltaTwo$
for the two correlations. A point is kept if $I(c)\le\tau$, where $\tau=\Tau$ is
the $\Confidence$ quantile of a $\chi^2$ distribution with $\NMoments$ degrees of
freedom.

The threshold is a rule for ruling out. It is not the critical value of a test,
and it does not define a credible region: the two level moments share a state
stream, $V$ is diagonal, and $\delta^2$ rather than sampling noise fills most of
$V$. Moving $\tau$ or $\delta$ after seeing which points survive would turn a
filter into a fit, so both are fixed before any run. Appendix~\ref{app:menu}
reports, as a disclosed sensitivity analysis, how the image changes when the level
tolerance is tightened.

The search uses no emulator; every point is simulated. Wave~1 is a Latin
hypercube of $\WavePoints$ points in $C$. Waves~2 and~3 each draw $\WavePoints$
new points in the bounding box of the points kept by the previous wave, widened
by $\BoxInflationPct\%$ of its span and clipped to $C$. In the code, NROY is a
finite set: the union over the three waves of the sampled points with
$I(c)\le\tau$. It approximates the level set $\{c\in C:I(c)\le\tau\}$ without
enumerating it. Membership means that the targets cannot rule a point out. It
does not mean the point is true, and NROY is not a posterior over $c$. An empty
NROY would be a finding: the pair of kernel and targets would be rejected.

\subsection{The image}

The object we report is the image of the policy effect over the kept points.
Let $\mathcal S$ be a declared uniform sample of $\ThetaSample$ NROY points. The
code reports
\begin{equation}\label{eq:image}
  \Theta_{\hm}=\bigl[\min_{c\in\mathcal S}\theta_2(c),\;\max_{c\in\mathcal S}\theta_2(c)\bigr],
  \qquad
  \Theta_{\expset}=\bigl[\min_{c\in\mathcal S}\dspill(c),\;\max_{c\in\mathcal S}\dspill(c)\bigr],
\end{equation}
where $\dspill$ is the spillover contrast defined with Foil~2 in
Section~\ref{sec:foil2}. Each image is the range over the sample, the declared
empirical image. It is a set, not a distribution: the range says which values the
targets leave open and says nothing about how likely each one is.

The endpoints are extrema of noisy estimates, so the code evaluates each endpoint
candidate again on fresh replications (Appendix~\ref{app:recheck}).

\subsection{The diagnostic points}

The analysts of Section~\ref{sec:analysts} work at three diagnostic points, which
we write $c_0$. Two belong to the calibrator, who fits the model and works at one
best-fitting point: the \emph{min-I point}, the NROY point with the smallest
implausibility, and the \emph{median-I point}, the point at the middle of NROY in
the order of $I$. Both are selected by $I$ alone, never by $\theta_2$. The third
is $\cstar$ itself, the real village, which belongs to the field experimenter.
We use three points rather than all $\ThetaSample$: running an analyst at every
kept point would hand her the calibration set, and default practice does not do
that.

Figure~\ref{fig:pipe} summarizes the order of operations. Every exogenous
quantity is fixed before any wave runs. If $\theta_2$ decided which points to
keep, the image would be a fit rather than a report.

\begin{figure}[t]
\centering
\begin{tikzpicture}[ecn, node distance=6mm and 6mm]
  \node[ecn box, text width=120mm] (freeze)
    {1~~fix the kernel, $C$, the targets $z$, $\delta$, $\tau$, the estimands and the rule for the points $c_0$};
  \node[ecn box, text width=120mm, below=of freeze] (hm)
    {2~~history matching: $I(c)\le\tau$ over three waves $\;\to\;$ NROY\\[0.15em]
     {\scriptsize not a posterior; an empty NROY is a finding}};
  \node[ecn box, text width=56mm, below=of hm, xshift=-32mm] (A)
    {\textbf{3~~image}\\[0.15em]
     $\Theta_{\hm}$ and $\Theta_{\expset}$ over $\ThetaSample$ sampled NROY points;
     endpoint recheck};
  \node[ecn box, text width=56mm, below=of hm, xshift=32mm] (B)
    {\textbf{4~~analysts}\\[0.15em]
     one panel at $c_0\in\{$min-I, median-I, $\cstar\}$;
     $\TradeoffR$ fresh panels for A0--A2};
  \node[ecn mark, text width=120mm, below=11mm of hm |- A.south] (ov)
    {each interval against the truth at $c_0$, the truth at $\cstar$, and its own image};
  \draw[ecn arrow] (freeze) -- (hm);
  \draw[ecn arrow] (hm.south) -- ++(0,-0.25) -| (A.north);
  \draw[ecn arrow] (hm.south) -- ++(0,-0.25) -| (B.north);
  \draw[ecn arrow] (A.south) -- ++(0,-0.3) -| (ov.north);
  \draw[ecn arrow] (B.south) -- ++(0,-0.3) -| (ov.north);
\end{tikzpicture}
\caption{Order of operations. The image and the analysts are two separate reports;
the evaluation compares them and never merges them into one object.}
\label{fig:pipe}
\end{figure}
